% GENERATED FILE. Edit canonical lessons, then run npm run generate:books.
% canonical-source-hash: fnv1a64:8cc636beb6dfce61
% canonical-lessons: ZH-C282-gai, ZH-W282-gai, ZH-C282-yinggai, ZH-W282-ying, ZH-C282-yinggai-grammar, ZH-C282-zuihao, ZH-C282-dei, ZH-C282-practice

\chapter{You Should Rest}
\label{ch:zh-you-should-rest}
\hlchaptermodality{282}

\begin{chapteropening}
\textbf{By the end of this chapter:} \emph{I can say what someone should or had better do.}

Giving advice with \zh{应该} and \zh{不应该}, saying it is time with \zh{该}, what is wise with \zh{最好} and what cannot be helped with děi, built on two new written characters: \zh{该} and \zh{应}.
\end{chapteropening}

\section[\texorpdfstring{\zh{该} (gāi)}{该 (gāi)}]{\zh{该} (gāi) — it's time to}

\label{lesson:ZH-C282-gai}



\begin{quote}
A new chapter, and it opens on the moment when something has to happen.

\begin{itemize}
\raggedright
  \item \emph{Recall:} say \zh{因为我喜欢中国} \zh{所以学汉语} — \emph{yīnwèi wǒ xǐhuan Zhōngguó, suǒyǐ xué Hànyǔ} — \textbf{R1}, the chapter you have finished
  \item \emph{Recall:} say \emph{lìzi}, an example — \textbf{R1}
  \item \emph{Recall:} say \emph{xiūxi}, to rest — \textbf{R3}, durable
\end{itemize}
\end{quote}

\subsection*{You'll want to know: \zh{该}}
\textbf{\zh{该}} — \emph{gāi} — \textquotedblleft{}it's time to; should (by now)\textquotedblright{}.

\begin{quote}
\textbf{\zh{我该回家了}} \emph{wǒ gāi huí jiā le} — It's time I went home.
\end{quote}

\begin{quote}
\textbf{\zh{该你了}} \emph{gāi nǐ le} — Your turn.
\end{quote}

The next lesson puts it in your hand.

\begin{cousinweb}
\textbf{\zh{该}} (\emph{gāi}) is inherited from Middle Chinese \emph{koj} (Guangyun, in Baxter's transcription), glossed 'complete' and 'all'. 'Should' and 'it's time to' are not given there, so how the word came by them is not settled.
\end{cousinweb}

\begin{grammarlens}[title={What you've built}]
With \zh{了} at the end, \zh{该} says the moment has come: \zh{我该回家了}, \zh{该你了}.

It sits before the verb, after the person, where \zh{要} and \zh{想} sit.
\end{grammarlens}

\subsection*{Your turn}
\begin{itemize}
\raggedright
  \item \emph{Say it:} \emph{gāi}
  \item \emph{Say it:} \emph{gāi}, once more
  \item \emph{Say it:} \zh{我该回家了} — \emph{wǒ gāi huí jiā le} — It's time I went home.
  \item \emph{Say it:} \zh{该你了} — \emph{gāi nǐ le} — Your turn.
  \item \emph{Recall:} say \emph{xiūxi}, then \zh{我该} \emph{xiūxi} \zh{了}: time I rested
\end{itemize}

\begin{tcolorbox}[breakable,colback=teal!4,colframe=teal!35!black,title={Before you move on}]
What does \zh{该} mean? (\textquotedblleft{}It's time to; should\textquotedblright{}.) What closes \zh{我该回家了}? (\textbf{\zh{了}: the moment has come}.)
\end{tcolorbox}
\section[\texorpdfstring{\zh{该} (gāi)}{该 (gāi)}]{\zh{该} — the left of \zh{说}, then the right of \zh{孩}}

\label{lesson:ZH-W282-gai}



\begin{quote}
The word first, then the character inside it.

\begin{itemize}
\raggedright
  \item \emph{Recall:} say \emph{gāi} — \textbf{R1}, one lesson back, before it decays
\end{itemize}

\begin{quote}
\zh{说} \zh{孩}
\end{quote}

\emph{Read it:} each one aloud, with nothing said first — \emph{shuō}, \emph{hái}
\end{quote}

\subsection*{Script — the left of \zh{说}, then the right of \zh{孩}}
\hlblockfigure{\detokenize{figures/ZH-W282-gai-filmstrip.pdf}}{How \zh{该} is written, stroke by stroke}

\begin{quote}
\zh{该}
\end{quote}

\textbf{Two pieces, and you write both already.}

\textbf{The left of \zh{说}, two strokes}: a dot, then a short horizontal that turns down and flicks up to the right.

\textbf{The right of \zh{孩}, six strokes}: a dot at the top; a horizontal under it; a stroke that falls to the lower left and turns up to the right; a long left-falling stroke; a shorter left-falling stroke under it; and a dot at the lower right.

2 + 6 = eight.

\subsection*{Your turn}
\begin{itemize}
\raggedright
  \item \emph{Write it:} \zh{该}, all eight, with the model in front of you
  \item \emph{Cover:} write it again from memory, and count the pieces, not the strokes
\end{itemize}

\begin{tcolorbox}[breakable,colback=teal!4,colframe=teal!35!black,title={Before you move on}]
How many strokes in \zh{该}? (\textbf{Eight}.) Where have you written its left side? (\textbf{On \zh{说}}.) And its right side? (\textbf{On \zh{孩}}.)

Source: \href{https://raw.githubusercontent.com/chanind/hanzi-writer-data/68d10a4b21150cae5e1ebbd223eed289cf32d90c/data/\%E8\%AF\%A5.json}{Hanzi Writer Data \zh{该}, pinned stroke paths and medians}.
\end{tcolorbox}
\section[\texorpdfstring{\zh{应该} (yīnggāi)}{应该 (yīnggāi)}]{\zh{应该} (yīnggāi) — should; ought to}

\label{lesson:ZH-C282-yinggai}



\begin{quote}
Before the new one: the character you wrote last, in the sentence it belongs to.

\begin{itemize}
\raggedright
  \item \emph{Read it:} \zh{我该回家了}, aloud — \emph{wǒ gāi huí jiā le}
  \item \emph{Recall:} say \zh{医生} — \emph{yīshēng}, a doctor — \textbf{R3}, durable
\end{itemize}
\end{quote}

\subsection*{You'll want to know: \zh{应该}}
\textbf{\zh{应该}} — \emph{yīnggāi} — \textquotedblleft{}should; ought to\textquotedblright{}.

\begin{quote}
\textbf{\zh{你应该去看医生}} \emph{nǐ yīnggāi qù kàn yīshēng} — You should see a doctor.
\end{quote}

The next lesson puts the first character in your hand. The second is the \zh{该} you wrote last.

\begin{cousinweb}
\textbf{\zh{应该}} (\emph{yīnggāi}) is built from \zh{应} and \zh{该}, both inherited from Middle Chinese. The Guangyun lists \zh{应} under two readings: \emph{'ing}, glossed with a word that can mean 'ought to', and \emph{'ingH}, glossed 'things answering each other'; yīng is the first. \zh{该} is \emph{koj}, glossed 'complete' (Guangyun, in Baxter's transcription).
\end{cousinweb}

\begin{grammarlens}[title={What you've built}]
\zh{该} says the moment has come; \zh{应该} gives advice, or says what is right.

Both sit before the verb: \textbf{\zh{你应该去看医生}} \emph{nǐ yīnggāi qù kàn yīshēng}. \zh{看医生} is to see a doctor.
\end{grammarlens}

\subsection*{Your turn}
\begin{itemize}
\raggedright
  \item \emph{Say it:} \emph{yīnggāi}
  \item \emph{Say it:} \emph{yīnggāi}, once more
  \item \emph{Say it:} \zh{你应该去看医生} — \emph{nǐ yīnggāi qù kàn yīshēng} — You should see a doctor.
  \item \emph{Recall:} say \zh{该}, then \zh{应该}: it's time, then you should
\end{itemize}

\begin{tcolorbox}[breakable,colback=teal!4,colframe=teal!35!black,title={Before you move on}]
What does \zh{应该} mean? (\textquotedblleft{}Should; ought to\textquotedblright{}.) How is it different from \zh{该}? (\textbf{\zh{该} says it's time; \zh{应该} gives advice}.)
\end{tcolorbox}
\section[\texorpdfstring{\zh{应} (yīng)}{应 (yīng)}]{\zh{应} — the roof of \zh{店}, two dots, a falling stroke and a base}

\label{lesson:ZH-W282-ying}



\begin{quote}
The word first, then the character inside it.

\begin{itemize}
\raggedright
  \item \emph{Recall:} say \emph{yīnggāi} — \textbf{R1}, one lesson back, before it decays
\end{itemize}

\begin{quote}
\zh{店} \zh{点}
\end{quote}

\emph{Read it:} each one aloud, with nothing said first — \emph{diàn}, \emph{diǎn}
\end{quote}

\subsection*{Script — the roof of \zh{店}, two dots, a falling stroke and a base}
\hlblockfigure{\detokenize{figures/ZH-W282-ying-filmstrip.pdf}}{How \zh{应} is written, stroke by stroke}

\begin{quote}
\zh{应}
\end{quote}

\textbf{Two pieces, and you write the first already.}

\textbf{The roof of \zh{店}, three strokes}: the dot on top, the bar, and a long left-falling stroke down the left side.

\textbf{Under it, four strokes}: a dot on the left, a dot in the middle, a left-falling stroke on the right, and last a long horizontal along the foot.

3 + 4 = seven.

\subsection*{Your turn}
\begin{itemize}
\raggedright
  \item \emph{Write it:} \zh{应}, all seven, with the model in front of you
  \item \emph{Cover:} write it again from memory, and count the pieces, not the strokes
\end{itemize}

\begin{tcolorbox}[breakable,colback=teal!4,colframe=teal!35!black,title={Before you move on}]
How many strokes in \zh{应}? (\textbf{Seven}.) Which character you write has the same roof? (\textbf{\zh{店}}.) Which stroke comes last? (\textbf{The long horizontal along the foot}.)

Source: \href{https://raw.githubusercontent.com/chanind/hanzi-writer-data/68d10a4b21150cae5e1ebbd223eed289cf32d90c/data/\%E5\%BA\%94.json}{Hanzi Writer Data \zh{应}, pinned stroke paths and medians}.
\end{tcolorbox}
\section[\texorpdfstring{\zh{你应该}… (nǐ yīnggāi …)}{你应该… (nǐ yīnggāi …)}]{\zh{你应该}… — you should}

\label{lesson:ZH-C282-yinggai-grammar}



\begin{quote}
A word you have, and the advice it gives.

\begin{itemize}
\raggedright
  \item \emph{Read it:} \zh{应该}, aloud — \emph{yīnggāi}
  \item \emph{Recall:} say \zh{我该回家了} — \emph{wǒ gāi huí jiā le} — \textbf{R2}
  \item \emph{Recall:} say \emph{zěnme bàn}, what should I do — \textbf{R3}, durable
\end{itemize}
\end{quote}

\begin{grammarlens}[title={Grammar Lens: \zh{应该} and \zh{不应该}, advice and what is right}]
\begin{quote}
\textbf{\zh{应该} before a verb gives advice, or says what is right; \zh{不应该} says what is not.}
\end{quote}

\begin{quote}
\textbf{\zh{你应该去看医生}} \emph{nǐ yīnggāi qù kàn yīshēng} — You should see a doctor. \textbf{\zh{你不应该这么晚回家}} \emph{nǐ bù yīnggāi zhème wǎn huí jiā} — You shouldn't come home so late.
\end{quote}

\zh{不} goes in front of \zh{应该}, and stays \emph{bù} before the level yīng.

To ask for advice: \emph{wǒ yīnggāi zěnme bàn?}, what should I do?

\zh{该} says the time has come; \zh{应该} says what is right, whatever the time.
\end{grammarlens}

\subsection*{Your turn}
Say these aloud:

\begin{itemize}
\raggedright
  \item \zh{你应该去看医生} — \emph{nǐ yīnggāi qù kàn yīshēng} — You should see a doctor.
  \item \zh{你不应该这么晚回家} — \emph{nǐ bù yīnggāi zhème wǎn huí jiā} — You shouldn't come home so late.
  \item \emph{wǒ yīnggāi zěnme bàn} — What should I do?
  \item one piece of advice for a friend who is ill
\end{itemize}

\begin{tcolorbox}[breakable,colback=teal!4,colframe=teal!35!black,title={Before you move on}]
Where does \zh{应该} go? (\textbf{Before the verb}.) How do you say someone shouldn't? (\textbf{\zh{不应该}, said \emph{bù yīnggāi}}.)
\end{tcolorbox}
\section[\texorpdfstring{\zh{最好} (zuìhǎo)}{最好 (zuìhǎo)}]{\zh{最好} (zuìhǎo) — had better}

\label{lesson:ZH-C282-zuihao}



\begin{quote}
Before the new one: advice, and the character it starts with.

Read these aloud:

\begin{itemize}
\raggedright
  \item \zh{你不应该这么晚回家}, aloud — \emph{nǐ bù yīnggāi zhème wǎn huí jiā}
  \item \zh{最后}, aloud — \emph{zuìhòu}
\end{itemize}
\end{quote}

\subsection*{You'll want to know: \zh{最好}}
\textbf{\zh{最好}} — \emph{zuìhǎo} — \textquotedblleft{}had better; it would be best to\textquotedblright{}.

\begin{quote}
\textbf{\zh{你最好早点儿回家}} \emph{nǐ zuìhǎo zǎo diǎnr huí jiā} — You'd better go home early.
\end{quote}

\begin{quote}
\emph{nǐ zuìhǎo bié qù} — You'd better not go.
\end{quote}

Every character in it is one you write.

\begin{cousinweb}
\textbf{\zh{最好}} (\emph{zuìhǎo}) is built from \zh{最} 'most' and \zh{好} 'good', both inherited from Middle Chinese, \emph{tswajH} and \emph{xawX} (Guangyun, in Baxter's transcription), glossed 'utmost' and 'good'. 'Best' is the plain sum of the two; 'had better' is the word's own use.
\end{cousinweb}

\begin{grammarlens}[title={What you've built}]
\zh{应该} says what is right; \zh{最好} says what is wise, the best of the choices.

It goes where \zh{应该} goes, before the verb: \textbf{\zh{你最好早点儿回家}} \emph{nǐ zuìhǎo zǎo diǎnr huí jiā}.
\end{grammarlens}

\subsection*{Your turn}
\begin{itemize}
\raggedright
  \item \emph{Say it:} \emph{zuìhǎo}
  \item \emph{Say it:} \emph{zuìhǎo}, once more
  \item \emph{Say it:} \zh{你最好早点儿回家} — \emph{nǐ zuìhǎo zǎo diǎnr huí jiā} — You'd better go home early.
  \item \emph{Say it:} \emph{nǐ zuìhǎo bié qù} — You'd better not go.
  \item \emph{Recall:} say \zh{最后}, then \zh{最好}: the same \zh{最}, two words
\end{itemize}

\begin{tcolorbox}[breakable,colback=teal!4,colframe=teal!35!black,title={Before you move on}]
What does \zh{最好} mean? (\textquotedblleft{}Had better\textquotedblright{}.) Where does it go? (\textbf{Before the verb, where \zh{应该} goes}.)
\end{tcolorbox}
\section[děi]{děi — to have to}

\label{lesson:ZH-C282-dei}



\begin{quote}
Before the new one: last lesson's word, and a character you know in another job.

\begin{itemize}
\raggedright
  \item \emph{Recall:} say \zh{最好} — \emph{zuìhǎo}, had better — \textbf{R1}, one lesson back
  \item \emph{Read it:} \zh{你说得很好}, aloud — \emph{nǐ shuō de hěn hǎo}
\end{itemize}
\end{quote}

\subsection*{You'll want to know: děi}
\emph{děi} — \textquotedblleft{}to have to; must\textquotedblright{}.

\begin{quote}
\textbf{\zh{我得回家了}} \emph{wǒ děi huí jiā le} — I have to go home now.
\end{quote}

\begin{quote}
\textbf{\zh{明天我得去上课}} \emph{míngtiān wǒ děi qù shàngkè} — I have to go to class tomorrow.
\end{quote}

It is written \zh{得}, the character of \zh{说得好}, but it is a different word with its own sound, \emph{děi}.

\begin{cousinweb}
\emph{děi} is written with \zh{得}, which is inherited from Middle Chinese \emph{tok} (Guangyun, in Baxter's transcription), glossed 'to get'. \emph{děi} is not the regular modern sound of \emph{tok}, and 'to have to' is not given there; how the character came to write this word is not settled.
\end{cousinweb}

\begin{grammarlens}[title={What you've built}]
\zh{应该} is advice, \zh{最好} is what is wise, and \emph{děi} leaves no choice: \zh{我得回家了}.

So \zh{得} in a sentence has two sounds: \emph{de} after a verb, as in \zh{说得好}, and \emph{děi} before one, as here.
\end{grammarlens}

\subsection*{Your turn}
\begin{itemize}
\raggedright
  \item \emph{Say it:} \emph{děi}
  \item \emph{Say it:} \emph{děi}, once more
  \item \emph{Say it:} \zh{我得回家了} — \emph{wǒ děi huí jiā le} — I have to go home now.
  \item \emph{Say it:} \zh{明天我得去上课} — \emph{míngtiān wǒ děi qù shàngkè} — I have to go to class tomorrow.
  \item \emph{Recall:} say \zh{你说得很好}, then \zh{我得回家了}: \emph{de}, then \emph{děi}
\end{itemize}

\begin{tcolorbox}[breakable,colback=teal!4,colframe=teal!35!black,title={Before you move on}]
What does \emph{děi} mean? (\textquotedblleft{}To have to\textquotedblright{}.) How is \zh{得} said in \zh{说得好}? (\textbf{de, light}.) And in \zh{我得回家了}? (\textbf{děi}.)
\end{tcolorbox}
\section[Practice]{\zh{你病了} \zh{应该去看医生}}

\label{lesson:ZH-C282-practice}



\begin{quote}
The chapter closes here, with a friend who is not well and wants to go to class.

\begin{itemize}
\raggedright
  \item \emph{Recall:} say \zh{我得回家了} — \emph{wǒ děi huí jiā le} — \textbf{R1}, one lesson back
  \item \emph{Recall:} say \zh{最好} — \emph{zuìhǎo} — \textbf{R2}
  \item \emph{Recall:} say \emph{hǎoxiàng}, as if — \textbf{R3}
\end{itemize}
\end{quote}

\subsection*{The exchange}
\begin{quote}
\textbf{A:} \emph{wǒ hǎoxiàng bìng le} — I think I'm ill.
\end{quote}

\begin{quote}
\textbf{B:} \zh{你应该去看医生}
\end{quote}

\begin{quote}
\emph{nǐ yīnggāi qù kàn yīshēng} — You should see a doctor.
\end{quote}

\begin{quote}
\textbf{A:} \zh{我得去上课}
\end{quote}

\begin{quote}
\emph{wǒ děi qù shàngkè} — I have to go to class.
\end{quote}

\begin{quote}
\textbf{B:} \zh{你病了} \zh{不应该去上课}
\end{quote}

\begin{quote}
\emph{nǐ bìng le, bù yīnggāi qù shàngkè} — You're ill. You shouldn't go to class.
\end{quote}

\begin{quote}
\textbf{A:} \zh{那我该做什么}
\end{quote}

\begin{quote}
\emph{nà wǒ gāi zuò shénme} — Then what should I do?
\end{quote}

\begin{quote}
\textbf{B:} \zh{你最好在家}
\end{quote}

\begin{quote}
\emph{nǐ zuìhǎo zài jiā} — You'd better stay at home.
\end{quote}

\begin{quote}
\textbf{A:} \zh{好} \zh{我该回家了}
\end{quote}

\begin{quote}
\emph{hǎo, wǒ gāi huí jiā le} — All right. Time I went home.
\end{quote}

Four words of advice and duty, in rising order of force: \zh{最好} is wise, \zh{应该} is right, \zh{该} says the time has come, \emph{děi} allows no choice.

\begin{tcolorbox}[breakable,skin=enhanced,colback=orange!6,colframe=orange!45!black,boxrule=0.5pt,arc=1mm,left=6pt,right=6pt,top=4pt,bottom=4pt,fonttitle=\bfseries,title={Read this}]
\emph{Read it:} the exchange through once without stopping, then answer the question below
\end{tcolorbox}

\begin{grammarlens}[title={What you've built — advice and duty}]
Four words before the verb, each with its own force:

\begin{itemize}
\raggedright
  \item \textbf{\zh{你应该去看医生}} \emph{nǐ yīnggāi qù kàn yīshēng} — what is right.
  \item \textbf{\zh{你最好在家}} \emph{nǐ zuìhǎo zài jiā} — what is wise.
  \item \textbf{\zh{我该回家了}} \emph{wǒ gāi huí jiā le} — the time has come.
  \item \textbf{\zh{我得去上课}} \emph{wǒ děi qù shàngkè} — no choice.
\end{itemize}

\textbf{Listen.} \emph{Nǐ juéde wǒ yīnggāi qù ma?}

\textbf{Speak.} Give a friend two pieces of advice about the weekend, one with \zh{应该} and one with \zh{最好}.

\emph{Write it:} \zh{应该} — the roof of \zh{店} first, then the right of \zh{孩} after the left of \zh{说}
\end{grammarlens}

\subsection*{Your turn}
\begin{itemize}
\raggedright
  \item \emph{Say it:} \zh{你病了} \zh{应该在家} — \emph{nǐ bìng le, yīnggāi zài jiā} — You're ill; you should stay at home.
  \item \emph{Say it:} \zh{你最好早点儿回家} — \emph{nǐ zuìhǎo zǎo diǎnr huí jiā} — You'd better go home early.
  \item \emph{Say it:} the whole exchange, both parts
  \item \emph{Write it:} \zh{你应该去看医生}
\end{itemize}

\begin{tcolorbox}[breakable,colback=teal!4,colframe=teal!35!black,title={Before you move on}]
Say you should see a doctor. (\textbf{\zh{你应该去看医生}}.) Say you'd better stay at home. (\textbf{\zh{你最好在家}}.) Say it's time you went home. (\textbf{\zh{我该回家了}}.)
\end{tcolorbox}
